\clearpage\begingroup
\HistoricalHead{193}{R}{WECHSELWIRKUNGEN ZWISCHEN GALVANISCHEN STRÖMEN UND MAGNETISMUS.}
\begin{center}
\textbf{Zweiter Abschnitt.}\\[8pt]
\textls[100]{Allgemeine Ausdrücke für die Kraft, mit welcher ein wie immer}\\
\textls[100]{geformter Leiter eines galvanischen Stroms auf ein magnetisches}\\
\textls[100]{Element in jedem Punkte des Raumes wirkt.}\\[10pt]
11.
\end{center}
\noindent\hspace*{\parindent}Wir bezeichnen die Länge des Stromleiters, von einem beliebigen Anfangspunkte in demselben bis zu einem unbestimmten Punkte desselben $Q$, in der Richtung des positiven Stromes gezählt, durch $s$; die Coordinaten in Beziehung auf drei beliebige sich rechtwinklicht schneidende Axen für $Q$ durch $x,y,z$; die bewegende Kraft, welche das Stromelement $QQ'=ds$ in einem Punkte des Raumes $P$, dessen Coordinaten $a,b,c$, ausübt, durch $\omega$, indem wir die Intensität des Stromes und das magnetische Element, auf welches gewirkt wird, $=1$ setzen; die drei partiellen mit den Coordinatenaxen parallelen Kräfte, in welche sich $\omega$ zerlegt, durch $\xi,\eta,\zeta$, so dass $\frac{\xi}{\omega},\frac{\eta}{\omega},\frac{\zeta}{\omega}$ die Cosinus der Winkel sind, welche die Richtung der Kraft $\omega$ mit den Coordinatenaxen machen; endlich $PQ$ durch $r$, und $PS$ drücke die Richtung der Kraft aus.

Die hier erforderlichen Entwicklungen werden sich am leichtesten übersehen lassen, wenn wir alle hier vorkommenden Richtungen durch Punkte auf einer willkührlichen Kugelfläche bezeichnen (vergl. \textit{Disquiss. genn. circa superficies curvas}, art. 1[*)]\begingroup\renewcommand{\thefootnote}{}\footnotetext{[*) Werke IV, S. 219.]}\endgroup). Es bezeichnen also (1), (2), (3) die Richtungen der drei Coordinatenaxen; (4) die Richtung der geraden Linie von $P$ nach $Q$; (5) die Richtung des Elements $QQ'$; (6) die Richtung der Kraft $\omega$, mit welcher das positive magnetische Fluidum 1 in $P$ durch das Stromelement $QQ'$ sollicitirt wird.
\HistoricalFoot{XI\textsubscript{1}.}{25}
\clearpage\SourceLeaf{199}\HistoricalHead{194}{L}{PHYSIK.\quad NACHLASS.}
\noindent\hspace*{\parindent}Nach dem, was oben gesagt ist, wird also $\omega=\frac{\sin(4)(5)}{rr}ds$ sein, ferner (6) ein Pol des grössten Kreises durch (4) und (5), und zwar derjenige, welcher von diesem Kreisbogen links liegt, wenn er in dem Sinne von (4) nach (5) durchlaufen wird, das Innere der Kugel als unten betrachtet. Die drei Axen der Coordinaten $x,y,z$ nehmen wir in solcher Ordnung liegend an, wie drei Richtungen nach vorne, nach der rechten und nach oben, unter welcher Voraussetzung also (1), (2), (3) diejenigen Pole der grössten Kreisbögen (3)(2), (1)(3), (2)(1) sein werden, welche gegen dieselben ebenso liegen wie (6) gegen (4)(5).

Aus dem allgemeinen Lehrsatz V (\textit{Disqu. Gen.} art. 2[*)]\begingroup\renewcommand{\thefootnote}{}\footnotetext{[*) Werke IV, S. 220.]}\endgroup) folgt nun, dass
\[
\begin{aligned}
&\cos(3)(4)\mathbin{.}\cos(2)(5)-\cos(2)(4)\mathbin{.}\cos(3)(5)\\
&\qquad=\sin(3)(2)\mathbin{.}\sin(4)(5)\mathbin{.}\cos(1)(6),
\end{aligned}
\]
\noindent oder da $\sin(3)(2)=1$,
\[
\cos(3)(4)\mathbin{.}\cos(2)(5)-\cos(2)(4)\mathbin{.}\cos(3)(5)=\frac{rr\omega}{ds}\mathbin{.}\cos(1)(6),
\]
\noindent oder da
\[
\begin{array}{ll}
\cos(3)(4)=\dfrac{z-c}{r},&\cos(2)(5)=\dfrac{dy}{ds}\\[5pt]
\cos(2)(4)=\dfrac{y-b}{r},&\cos(4)(5)=\dfrac{dz}{ds}
\end{array}
\]
\noindent [folgt weiter:]
\[
\omega\mathbin{.}\cos(1)(6)=\xi,
\]
\[
\begin{array}{@{}l@{\hspace{6em}}l}
\mathrm{I.}&\xi=\dfrac{(z-c)dy-(y-b)dz}{r^3}
\end{array}
\]
\noindent und ebenso erhält man
\[
\begin{array}{@{}l@{\hspace{6em}}l}
\mathrm{II.}&\eta=\dfrac{(x-a)dz-(z-c)dx}{r^3}\\[8pt]
\mathrm{III.}&\zeta=\dfrac{(y-b)dx-(x-a)dy}{r^3}.
\end{array}
\]
\noindent Man würde für $\xi,\eta,\zeta$ dieselben Ausdrücke mit entgegengesetzten Zeichen erhalten, wenn man für die drei Coordinatenaxen die entgegengesetzte Ordnung gewählt hätte.
\begin{center}\rule{.18\textwidth}{.35pt}\end{center}
\clearpage\endgroup
